% Zirui (Ray) Chen — CV
% Single source of truth for assets/pdf/cv_ziruichen.pdf. Build with `make`.
% Reproduces the look of the previous Word-exported CV: Helvetica body,
% ruled small-caps-free section heads, right-aligned years, serif running header.

\documentclass[10pt,letterpaper]{article}

\usepackage[T1]{fontenc}
\usepackage[utf8]{inputenc}
\usepackage{textcomp}

% Helvetica (URW Nimbus Sans) as the default family, matching the original export.
\usepackage{helvet}
\renewcommand{\familydefault}{\sfdefault}

\usepackage[margin=0.85in,top=0.9in,bottom=0.8in]{geometry}
\usepackage{titlesec}
\usepackage{fancyhdr}
\usepackage[hidelinks]{hyperref}

\setlength{\parindent}{0pt}
\setlength{\parskip}{0pt}
\pagenumbering{arabic}

% Match the Word original's typography: ragged right, no hyphenation, single
% sentence spacing. Without these, LaTeX justifies and hyphenates by default,
% which reads visibly differently from the previous export.
\raggedright
\frenchspacing
\hyphenpenalty=10000
\exhyphenpenalty=10000
\sloppy

% Keep a citation from splitting across the page break.
\widowpenalty=10000
\clubpenalty=10000

% --- Running header: serif, as in the original (page 1 is bare) -----------------
\pagestyle{fancy}
\fancyhf{}
\renewcommand{\headrulewidth}{0pt}
\fancyhead[L]{\rmfamily\footnotesize Zirui Chen}
\fancyhead[R]{\rmfamily\footnotesize\thepage}

% --- Section heads: regular weight, ruled underneath ---------------------------
\titleformat{\section}
  {\normalfont\large}{}{0pt}{\MakeUppercase}
  [\vspace{-8pt}\rule{\textwidth}{0.6pt}]
\titlespacing*{\section}{0pt}{14pt}{6pt}

% --- Helpers -------------------------------------------------------------------
% Institution / role line with a right-aligned year.
\newcommand{\entry}[2]{\textbf{#1}\hfill #2\par}
% Citation with a hanging indent, as in the original.
\newcommand{\cit}[1]{\begingroup\hangindent=0.32in\hangafter=1 #1\par\endgroup\vspace{3pt}}
% Run-in bold label for the coursework / skills blocks.
\newcommand{\runin}[1]{\textbf{#1}\ }

\begin{document}
\thispagestyle{empty}

% --- Masthead ------------------------------------------------------------------
{\fontsize{22}{26}\selectfont Zirui (Ray) Chen\par}
\vspace{6pt}
Krieger Hall 107, 3400 N Charles St.\hfill Email: \href{mailto:zchen160@jh.edu}{zchen160@jh.edu}\par
Baltimore, MD 21218\par

% --- Education -----------------------------------------------------------------
\section{Education}
\entry{Johns Hopkins University}{2023-present}
Ph.D. Cognitive Science\par
Advisor: Michael F. Bonner\par
\vspace{8pt}

\entry{Johns Hopkins University}{2022-2023}
M.A. Cognitive Science\par
Advisor: Michael F. Bonner\par
\vspace{8pt}

\entry{Emory University}{2018-2022}
B.S. Neuroscience \& Behavioral Biology, B.A. Computer Science (GPA=3.94)\par
\textit{Summa cum laude}\par
Honors thesis: The global color of a scene is relevant for human visual scene discrimination.\par
Advisor: Daniel D. Dilks\par

% --- Publications --------------------------------------------------------------
\section{Publications}
\cit{\textbf{Chen, Z.}, Isik, L., \& Bonner, M.F. (2026). \href{https://doi.org/10.64898/2026.04.27.720701}{Multidimensional dynamics of object representations in the human visual system}. \textit{bioRxiv}.}
\cit{\textbf{Chen, Z.}, \& Bonner, M.F. (2025). \href{https://doi.org/10.1126/sciadv.adw7697}{Universal dimensions of visual representation}. \textit{Science Advances}, 11, eadw7697.}
\cit{Cheng, A., \textbf{Chen, Z.}, \& Dilks, D. D. (2023). \href{https://doi.org/10.1016/j.neuroimage.2023.119935}{A stimulus-driven approach reveals vertical luminance gradient as a stimulus feature that drives human cortical scene selectivity}. \textit{NeuroImage}, 269, 119935.}

% --- Talks ---------------------------------------------------------------------
\section{Talk Presentations}
\cit{\textbf{Chen, Z.}, \& Bonner, M.F. (2023), Canonical Dimensions of Neural Visual Representation. \textit{Annual Meeting of the Vision Sciences Society 2023}.}
\cit{\textbf{Chen, Z.}, Cheng, A., \& Dilks, D.D. (2020), Characterizing an image property relevant for human cortical scene processing. \textit{Summer Undergraduate Research Experience (SURE) Virtual Symposium 2020}.}

% --- Posters -------------------------------------------------------------------
\section{Poster Presentations}
\cit{\textbf{Chen, Z.}, \& Bonner, M.F. (2026), Revealing core representational axes of neural networks through learning dynamics and brain alignment. \textit{Conference on Cognitive Computational Neuroscience 2026}.}
\cit{\textbf{Chen, Z.}, \& Bonner, M.F. (2026), Learning dynamics reveal the natural basis of visual representation. \textit{Annual Meeting of the Vision Sciences Society 2026}.}
\cit{\textbf{Chen, Z.}, Isik, L., \& Bonner, M.F. (2025), Neural Dimensionality and Temporal Dynamics of Visual Representations. \textit{Conference on Cognitive Computational Neuroscience 2025}.}
\cit{\textbf{Chen, Z.}, Isik, L., \& Bonner, M.F. (2025), Unraveling the geometry of neural representational dynamics in rapid visual processing. \textit{Annual Meeting of the Vision Sciences Society 2025}.}
\cit{\textbf{Chen, Z.}, \& Bonner, M.F. (2024), Accounting for the reliability of deep neural networks in representational modeling. \textit{Conference on Cognitive Computational Neuroscience 2024}.}
\cit{\textbf{Chen, Z.}, \& Bonner, M.F. (2024), How to estimate noise ceilings for computational models of visual cortex. \textit{Annual Meeting of the Vision Sciences Society 2024}.}
\cit{\textbf{Chen, Z.}, \& Bonner, M.F. (2023), Canonical Dimensions of Vision. \textit{Conference of Cognitive Computational Neuroscience 2023}.}
\cit{Elmoznino, E., \textbf{Chen, Z.}, \& Bonner, M.F. (2022), Latent dimensionality scales with the performance of deep learning models of visual cortex. \textit{Conference of Cognitive Computational Neuroscience 2022}.}
\cit{\textbf{Chen, Z.}, Cheng, A., \& Dilks, D.D. (2021), Uncovering the visual features relevant to human visual scene discrimination. \textit{Summer Undergraduate Research Experience (SURE) Virtual Symposium 2021}.}

% --- Coursework ----------------------------------------------------------------
\section{Relevant Coursework}
\runin{Computer Science:} Analysis of Algorithms; Artificial Intelligence; Computational Linguistics; Data Structure and Algorithms; Machine Learning.\par
\vspace{4pt}
\runin{Mathematics:} Differential Equations; Linear Algebra; Linear Optimization; Nonlinear Optimization; Multivariable Calculus; Numerical Analysis; Real Analysis.\par
\vspace{4pt}
\runin{Psychology/Neuroscience/Cognitive Science:} Bayesian Inference; Behavioral Neuroscience; Brain Imaging; Cognition; Foundations of Cognitive Science; Neurobiology; Neuroeconomics: Decision Making; Perspectives in Neuroscience \& Behavior; Predictive Modeling; Theory and Modeling of Information Coding in Neural Activity.\par

% --- Skills --------------------------------------------------------------------
\section{Skills}
\runin{Programming languages:} Python, MATLAB, R, Java, Bash, C, JavaScript, HTML/CSS, Swift.\par
\vspace{4pt}
\runin{Methods and tools:} Machine learning, deep learning, fMRI pre-processing (FSL, fMRIPrep), image processing, experiment interface design (Psychtoolbox, MTurk, Prolific).\par

% --- Teaching ------------------------------------------------------------------
\section{Teaching Experience}
\entry{Teaching Assistant, Johns Hopkins University}{2025}
Course: Introduction to Cognitive Neuropsychology\par
Professor: Michael McCloskey\par
\vspace{8pt}

\entry{Teaching Assistant, Johns Hopkins University}{2025}
Course: Neuroscience: Cognitive\par
Professor: Michael F. Bonner\par
\vspace{8pt}

\entry{Teaching Assistant, Johns Hopkins University}{2024}
Course: Deep Learning for Cognitive Neuroscience\par
Professor: Leyla Isik\par
\vspace{8pt}

\entry{Teaching Assistant, Johns Hopkins University}{2023}
Course: Probabilistic Models of the Visual Cortex\par
Professor: Alan Yuille\par
\vspace{8pt}

\entry{Teaching Assistant, Emory University}{2021}
Course: Non-Linear Optimization\par
Professor: Elizabeth Newman\par
\vspace{8pt}

\entry{Teaching Assistant, Emory University}{2021}
Course: Computer Architecture/Machine Level Programming\par
Professor: Shun Cheung\par

\end{document}
